%%%%%%%%%%%%%%%%%%%%%%%%%%%%%%%%%%%%%%%%%%%%%%%%%%%%%%%%%%%%%%%%
%
%  Template for homework of Machine Learning.
%%%%%%%%%%%%%%%%%%%%%%%%%%%%%%%%%%%%%%%%%%%%%%%%%%%%%%%%%%%%%%%%

\documentclass[11pt,letter,notitlepage]{ctexart}

\usepackage{fontspec}
\setCJKmainfont{KaiTi}[BoldFont=KaiTi,ItalicFont=KaiTi]

% Mise en page
\usepackage[left=2cm, right=2cm, lines=45, top=0.8in, bottom=0.7in]{geometry}
\usepackage{fancyhdr}
\usepackage{fancybox}
\usepackage{graphicx}
\usepackage{pdfpages}
\usepackage{enumitem}
\usepackage{algorithm}
\usepackage{algorithmic}
\usepackage{array}
\usepackage{booktabs}
\renewcommand{\headrulewidth}{1.5pt}
\renewcommand{\footrulewidth}{1.5pt}
\newcommand\Loadedframemethod{TikZ}
\usepackage[framemethod=\Loadedframemethod]{mdframed}
\usepackage{amssymb,amsmath}
\usepackage{amsthm}
\usepackage{thmtools}
\usepackage{bm}
\usepackage{multirow}
\usepackage{makecell}
\usepackage{tabularx}
\usepackage{longtable}
\usepackage{float}
\usepackage{caption}
\usepackage[edges]{forest}

\setlength{\topmargin}{0pt}
\setlength{\textheight}{9in}
\setlength{\headheight}{0pt}
\setlength{\oddsidemargin}{0.25in}
\setlength{\textwidth}{6in}

%%%%%%%%%%%%%%%%%%%%%%%%
%%%%%% Define math operator %%%%%
%%%%%%%%%%%%%%%%%%%%%%%%
\DeclareMathOperator*{\argmin}{\bf argmin}
\DeclareMathOperator*{\argmax}{\bf argmax}
\DeclareMathOperator*{\relint}{\bf relint\,}
\DeclareMathOperator*{\dom}{\bf dom\,}
\DeclareMathOperator*{\intp}{\bf int\,}

\DeclareMathOperator*{\cl}{\bf cl\,}
\DeclareMathOperator*{\bd}{\bf bd\,}
\DeclareMathOperator*{\conv}{\bf conv\,}
\DeclareMathOperator*{\epi}{\bf epi\,}

%%%%%%%%%%%%%%%%%%%%%%%
\pagestyle{fancy}

%%%%%%%%%%%%%%%%%%%%%%%%
%% Define the Exercise environment %%
%%%%%%%%%%%%%%%%%%%%%%%%
\mdtheorem[
topline=false,
rightline=false,
leftline=false,
bottomline=false,
leftmargin=-10,
rightmargin=-10
]{exercise}{\textbf{Exercise}}

%%%%%%%%%%%%%%%%%%%%%%%%
%% Define the Problem environment %%
%%%%%%%%%%%%%%%%%%%%%%%%
\mdtheorem[
topline=false,
rightline=false,
leftline=false,
bottomline=false,
leftmargin=-10,
rightmargin=-10
]{problem}{\textbf{Problem}}

%%%%%%%%%%%%%%%%%%%%%%%
%% Define the Solution Environment %%
%%%%%%%%%%%%%%%%%%%%%%%
\declaretheoremstyle
[
spaceabove=0pt,
spacebelow=0pt,
headfont=\normalfont\bfseries,
notefont=\mdseries,
notebraces={(}{)},
headpunct={:},
headindent={},
postheadspace={ },
bodyfont=\normalfont,
qed=$\blacksquare$,
preheadhook={\begin{mdframed}[style=myframedstyle]},
postfoothook=\end{mdframed},
]{mystyle}

\declaretheorem[style=mystyle,title=Solution]{solution}

\mdfdefinestyle{myframedstyle}{%
    topline=false,
    rightline=false,
    leftline=false,
    bottomline=false,
    skipabove=-6ex,
    leftmargin=-10,
    rightmargin=-10,
    backgroundcolor=black!5}

%%%%%%%%%%%%%%%%%%%%%%%
%% Useful commands
%%%%%%%%%%%%%%%%%%%%%%%
\renewcommand{\eqref}[1]{Eq.~(\ref{#1})}
\newcommand{\figref}[1]{Fig.~\ref{#1}}
\newcommand{\proj}[2]{\textbf{P}_{#2} (#1)}
\newcommand{\lspan}[1]{\textbf{span} (#1)}
\newcommand{\rank}[1]{\textbf{rank} (#1)}
\newcommand{\tr}{\operatorname{tr}}
\newcommand{\RNum}[1]{\uppercase\expandafter{\romannumeral #1\relax}}
\newcommand{\mb}[1]{\mathbf{#1}}

% Definition environment
\theoremstyle{definition}
\newtheorem{definition}{Definition}

%%%%%%%%%%%%%%%%%%%%%%%%%%%%%%%%%%%%%%%%%%
%% Homework info
%%%%%%%%%%%%%%%%%%%%%%%%%%%%%%%%%%%%%%%%%%
\newcommand{\lec}{\text{肖力}}
\newcommand{\posted}{\text{2026.3.26}}
\newcommand{\due}{\text{2026.4.5}}
\newcommand{\hwno}{\text{2}}

%%%%%%%%%%%%%%%%%%%%%%%%%%%%%%%%%%%%
%% Put your information here
%%%%%%%%%%%%%%%%%%%%%%%%%%%%%%%%%%%%
\newcommand{\name}{\text{陈璟皓}}
\newcommand{\id}{\text{PB23010359}}

\lhead{\textbf{\name}}
\rhead{\textbf{\id}}
\chead{\textbf{Homework \hwno}}

\begin{document}
    \vspace*{-4\baselineskip}
    \thispagestyle{empty}

    \begin{center}
        {\bf\large 机器学习B}
    \end{center}

    \noindent
    Lecturer: \lec
    \hfill
    Homework \hwno
    \\
    Posted: \posted
    \hfill
    Due: \due
    \\
    Name: \name
    \hfill
    ID: \id

    \noindent
    \rule{\textwidth}{2pt}

    \medskip

    %%%%%%%%%%%%%%%%%%%%%%%%%%%%%%%%%%%%%%%%%%%%%%%%%%%%%%%%%%%%
    % Exercise 1
    %%%%%%%%%%%%%%%%%%%%%%%%%%%%%%%%%%%%%%%%%%%%%%%%%%%%%%%%%%%%
    \begin{exercise}[平均绝对误差与拉普拉斯分布的极大似然推导]
    给出ch03-01中最后一页PPT的推导证明：

    以最小化平均绝对值误差的学习算法，那它等价于噪声误差服从拉普拉斯分布的极大似然估计。

    如果随机变量的概率密度函数分布为
    \[
    f(x\mid \mu,b)=\frac{1}{2b}\exp\left(-\frac{|x-\mu|}{b}\right)
    \]
    \[
    =\frac{1}{2b}
    \begin{cases}
    \exp\left(-\dfrac{\mu-x}{b}\right), & \text{if } x<\mu,\\[1em]
    \exp\left(-\dfrac{x-\mu}{b}\right), & \text{if } x\ge \mu,
    \end{cases}
    \]
    那么它就是拉普拉斯分布。其中，$\mu$ 是位置参数，$b>0$ 是尺度参数。

    \end{exercise}

    \newpage

    \begin{solution}

    \end{solution}

    \newpage

    %%%%%%%%%%%%%%%%%%%%%%%%%%%%%%%%%%%%%%%%%%%%%%%%%%%%%%%%%%%%
    % Exercise 2
    %%%%%%%%%%%%%%%%%%%%%%%%%%%%%%%%%%%%%%%%%%%%%%%%%%%%%%%%%%%%
    \begin{exercise}[线性判别分析计算题]
某教师希望根据学生的平时表现，预测其是否能够通过课程考核。

选取以下两个指标作为特征：

\begin{itemize}
    \item $x_1$：平时作业成绩
    \item $x_2$：课堂出勤率
\end{itemize}

分类结果：

\begin{itemize}
    \item 第1类：通过
    \item 第2类：未通过
\end{itemize}

\begin{center}
\begin{tabular}{>{\centering\arraybackslash}p{1.5cm}
                >{\centering\arraybackslash}p{3cm}
                >{\centering\arraybackslash}p{3cm}
                >{\centering\arraybackslash}p{2cm}}
\hline
学生 & 作业成绩 $x_1$ & 出勤率 $x_2$ & 类别 \\
\hline
1  & 85 & 90 & 通过 \\
2  & 78 & 85 & 通过 \\
3  & 92 & 88 & 通过 \\
4  & 74 & 80 & 通过 \\
5  & 88 & 95 & 通过 \\
6  & 60 & 65 & 未通过 \\
7  & 55 & 70 & 未通过 \\
8  & 62 & 60 & 未通过 \\
9  & 58 & 68 & 未通过 \\
10 & 65 & 72 & 未通过 \\
\hline
\end{tabular}
\end{center}

\vspace{1em}

\noindent
基于上面给出的 10 个训练样本，计算线性判别分析模型，并预测

假设某学生：
\begin{itemize}
    \item 作业成绩：75
    \item 出勤率：82
\end{itemize}

的分类结果。
    \end{exercise}

    \newpage

    \begin{solution}

    \end{solution}

    \newpage

    %%%%%%%%%%%%%%%%%%%%%%%%%%%%%%%%%%%%%%%%%%%%%%%%%%%%%%%%%%%%
    % Exercise 3
    %%%%%%%%%%%%%%%%%%%%%%%%%%%%%%%%%%%%%%%%%%%%%%%%%%%%%%%%%%%%
    \begin{exercise}[决策树]
画出课本图4.8中决策树的获得过程。
    \end{exercise}

\begin{figure}[htbp]
\centering
\renewcommand{\thefigure}{4.8}
\begin{forest}
for tree={
    draw,
    rounded corners,
    align=center,
    minimum width=2.6cm,
    minimum height=0.9cm,
    s sep=10mm,
    l sep=12mm,
    edge={thick},
    font=\small,
    parent anchor=south,
    child anchor=north,
    if n children=0{
        ellipse,
        minimum width=1.8cm,
        minimum height=0.8cm,
    }{},
}
[{纹理=?}
    [{密度$\leq 0.381$?}, edge label={node[midway, above left]{清晰}}
        [坏瓜, edge label={node[midway, left]{是}}]
        [好瓜, edge label={node[midway, right]{否}}]
    ]
    [{触感=?}, edge label={node[midway, above]{稍糊}}
        [好瓜, edge label={node[midway, left]{硬挺}}]
        [坏瓜, edge label={node[midway, right]{软粘}}]
    ]
    [坏瓜, edge label={node[midway, above right]{模糊}}]
]
\end{forest}
\caption{在西瓜数据集3.0上基于信息增益生成的决策树}
\end{figure}

\newpage

\begin{solution}

\end{solution}

\end{document}